% !PRINTED 534
\weberpage{534}

%<<CONT
the $a_{1,r}, a_{2,r}, \ldots, a_{r-1,r}$ can all be assumed equal to
zero.

If we denote by $\Delta$ the fundamental number of the field $\Omega$,
the discriminant of the system $\alpha_1, \alpha_2, \ldots, \alpha_n$
determined by (3) is obtained according to the formula \S.\,144, (11):
\weq{5}{
\Delta(\alpha_1, \alpha_2, \ldots, \alpha_n) = a_{1,1}^2\, a_{2,2}^2
\ldots a_{n,n}^2\, \Delta.
}
This is an integral rational number different from zero, and
consequently the quantities $\alpha$ are a basis of $\Omega$.

In order, then, to show that the system $\alpha_1, \alpha_2, \ldots,
\alpha_n$ so determined is a basis of $\mu$, it still remains to be
proved that every integer $\alpha$ divisible by $\mu$ can be
represented in the form (2). To carry out this proof, let us assume
any index $r \leqq n$, and seek the condition for an integer of the
form
\weq{6}{
\gamma_r = h_1\,\omega_1 + h_2\,\omega_2 + \cdots + h_r\,\omega_r
}
to be divisible by $\mu$, when $h_1, h_2, \ldots, h_r$ are integral
rational numbers.

First of all it follows that $h_r$ must be divisible by $a_{r,r}$.
For if we denote by $q_r$ the residue of the division of $h_r$ by
$a_{r,r}$, and set
\[
h_r = l_r\,a_{r,r} + q_r, \qquad 0 \leqq q_r < a_{r,r},
\]
then by (6) we obtain
\[
\gamma_r - l_r\,\alpha_r = (h_1 - l_r\,a_{1,r})\,\omega_1 +
(h_2 - l_r\,a_{2,r})\,\omega_2 + \cdots + q_r\,\omega_r,
\]
and this number would also have to be divisible by $\mu$. By the
definition of $a_{r,r}$, however, this is possible only if $q_r = 0$.
But then $\gamma_r - l_r\,\alpha_r$ receives the expression:
\weq{7}{
\gamma_r - l_r\,\alpha_r = h_1'\,\omega_1 + h_1'\,\omega_2 + \cdots +
h_{r-1}'\,\omega_{r-1}, % [sic]
}
and thus becomes of the same form as (6), only that $r - 1$ takes the
place of $r$; and in (7) the same mode of reasoning is to be repeated.
Accordingly, by complete induction we obtain the theorem:

\gesperrt{An integer of the field $\Omega$ representable in the form
$h_1\,\omega_1 + h_2\,\omega_2 + \cdots + h_r\,\omega_r$ is always
then and only then divisible by $\mu$, when it is representable in
the form}
\[
l_1\,\alpha_1 + l_2\,\alpha_2 + \cdots + l_r\,\alpha_r
\]
\gesperrt{in which the coefficients $l_1, l_2, \ldots, l_r$ are
integral rational numbers.}

If one sets $r = n$ in this theorem, one has the proof that the
system of numbers (3) is a basis of $\mu$.
